\section{ПРОЕКТИРОВАНИЕ ПРОГРАММНОГО КОМПЛЕКСА}

\subsection{Выбор стека технологий}

При проектировании программного комплекса были выбраны следующие технологии:

\begin{itemize}
    \item \textbf{Python 3.12} --- основной язык разработки, обеспечивающий высокую скорость прототипирования и обширную экосистему библиотек для научных вычислений;
    \item \textbf{Pygame 2.6} --- библиотека для создания графического интерфейса и визуализации, предоставляющая низкоуровневый доступ к 2D-рендерингу при минимальных зависимостях;
    \item \textbf{NumPy} --- библиотека для работы с многомерными массивами, используемая для представления сеточной карты мира;
    \item \textbf{Pandas} --- библиотека для обработки табличных данных, применяемая в подсистеме пакетного тестирования для агрегирования результатов экспериментов;
    \item \textbf{Matplotlib} --- библиотека для построения графиков, используемая для автоматической генерации визуализаций результатов пакетных экспериментов;
    \item \textbf{JSON} --- формат сериализации сценариев (карт), обеспечивающий совместимость и читаемость;
    \item \textbf{PyInstaller} --- инструмент для сборки Python-приложения в автономный исполняемый файл, обеспечивающий кроссплатформенное распространение без необходимости установки интерпретатора Python.
\end{itemize}

Выбор Python обусловлен тем, что целью работы является сравнение алгоритмов, а не достижение максимальной производительности отдельных реализаций. Python позволяет реализовать все алгоритмы в единой среде с одинаковыми накладными расходами интерпретатора, что обеспечивает корректность сравнения. Библиотека Pygame была выбрана вместо веб-фреймворков или Qt, поскольку она обеспечивает прямой контроль над кадровым циклом, что критично для покадровой визуализации работы алгоритмов. Использование PyInstaller позволяет собирать приложение в единый исполняемый файл для операционных систем Windows, Linux и macOS, что обеспечивает кроссплатформенность и упрощает распространение.

\subsection{Архитектура программного комплекса}

Архитектура программного комплекса построена по модульному принципу с чётким разделением ответственности между компонентами. На рисунке~\ref{fig:architecture} представлена общая схема архитектуры.

% === МЕСТО ДЛЯ РИСУНКА ===
\begin{figure}
  \centering
  %\fbox{\parbox{0.85\textwidth}{\centering\vspace{5cm}\textit{Диаграмма архитектуры программного комплекса}\vspace{5cm}}}
  \includegraphics[width=0.6\textwidth]{inc/architecture}
  \caption{Архитектура программного комплекса}
  \label{fig:architecture}
\end{figure}

Программный комплекс состоит из следующих модулей:

\begin{itemize}
    \item \textbf{core} --- базовые модели данных: мир (сеточная карта), агент, задача;
    \item \textbf{mapf} --- реализации MAPF-алгоритмов, объединённые общим базовым классом;
    \item \textbf{mrta} --- реализации MRTA-алгоритмов с единым интерфейсом;
    \item \textbf{simulation} --- симулятор, координирующий планирование и пошаговое выполнение;
    \item \textbf{batch} --- подсистема пакетного тестирования: конфигурация, генерация сценариев, запуск экспериментов, построение графиков;
    \item \textbf{ui} --- экраны графического интерфейса: главное меню, редактор карт, выбор алгоритма, экраны батч-режима;
    \item \textbf{visualization} --- отрисовка симуляции: карта, агенты, пути, панель метрик;
    \item \textbf{registry} --- центральный реестр алгоритмов;
    \item \textbf{scenario} --- модель сценария и сериализация в JSON.
\end{itemize}

Файловая структура проекта организована следующим образом:

\begin{verbatim}
src/
├── main.py                 # Точка входа
├── registry.py             # Реестр алгоритмов
├── scenario.py             # Модель сценария, JSON I/O
├── core/
│   ├── world.py            # Сеточная карта
│   ├── agent.py            # Агент
│   └── task.py             # Задача (MRTA)
├── mapf/
│   ├── planner.py          # Базовый класс MAPFPlanner
│   ├── a_star.py           # Time-expanded A*
│   ├── cbs/                # CBS (A* и SIPP)
│   ├── ecbs/               # ECBS
│   ├── eecbs/              # EECBS
│   ├── m_star/             # M*
│   ├── ca_star/            # Cooperative A*
│   ├── whca_star/          # WHCA*
│   ├── pibt/               # PIBT
│   └── sipp/               # SIPP (low-level)
├── mrta/
│   ├── planner.py          # Базовый класс MRTAPlanner
│   ├── hungarian.py        # Венгерский алгоритм
│   ├── greedy.py           # Жадный алгоритм
│   ├── sequential_auction.py
│   ├── min_cost_flow.py
│   └── combinatorial_auction.py
├── simulation/
│   └── simulator.py        # Симулятор
├── batch/
│   ├── config.py           # Конфигурация эксперимента
│   ├── runner.py           # Запуск экспериментов
│   ├── generator.py        # Генерация сценариев
│   └── plotter.py          # Построение графиков
├── ui/                     # Экраны интерфейса
└── visualization/          # Отрисовка симуляции
\end{verbatim}

\subsection{Базовые классы алгоритмов}

Ключевым архитектурным решением является использование паттерна <<Стратегия>> (Strategy) для подключения алгоритмов. Все MAPF- и MRTA-алгоритмы реализуют единый интерфейс, определённый базовыми классами \texttt{MAPFPlanner} и \texttt{MRTAPlanner}.

\subsubsection{Класс MAPFPlanner}

Базовый класс \texttt{MAPFPlanner} определяет контракт для всех алгоритмов поиска бесконфликтных маршрутов:

\begin{itemize}
    \item метод \texttt{plan(world, agents)} принимает модель мира и список агентов, возвращает словарь вида \texttt{\{agent\_id: [(pos, time), ...]\}} или \texttt{None} при неудаче;
    \item атрибуты класса содержат метаданные алгоритма: отображаемое имя (\texttt{DISPLAY\_NAME}), тип архитектуры (\texttt{IS\_CENTRALIZED}), режим работы (\texttt{IS\_ONLINE}).
\end{itemize}

Метаданные используются интерфейсом для отображения информации об алгоритме и симулятором для выбора стратегии выполнения. Единый формат возвращаемого значения --- список позиций с временными метками --- обеспечивает взаимозаменяемость алгоритмов: любой MAPF-планировщик может быть использован в качестве навигационного модуля для MRTA без модификации кода симулятора.

\subsubsection{Класс MRTAPlanner}

Базовый класс \texttt{MRTAPlanner} определяет контракт для алгоритмов распределения задач:

\begin{itemize}
    \item метод \texttt{plan(world, agents, tasks)} принимает мир, агентов и задачи, возвращает словарь \texttt{\{agent\_id: [Task, ...]\}};
    \item атрибут \texttt{ASSIGNMENT\_MODE} определяет режим назначения: \texttt{"online"} или \texttt{"queue"}.
\end{itemize}

Режим назначения является ключевым проектным решением, определяющим поведение симулятора:

\begin{itemize}
    \item в режиме \texttt{online} каждому агенту назначается по одной задаче; после её завершения планировщик вызывается повторно для освободившегося агента с оставшимися задачами;
    \item в режиме \texttt{queue} планировщик строит полную цепочку задач для каждого агента за один вызов; симулятор последовательно выполняет элементы цепочки без перепланирования.
\end{itemize}

Такое разделение позволяет реактивным алгоритмам (венгерский, жадный) адаптироваться к изменениям в реальном времени, а планирующим алгоритмам (аукционы, потоковые) оптимизировать глобальную стоимость назначения.

\subsection{Модель данных}

\subsubsection{Класс World}

Класс \texttt{World} представляет среду в виде двумерной сеточной карты. Внутреннее представление --- массив NumPy размером $H \times W$, где значение 1 обозначает препятствие. Метод \texttt{is\_free(x, y, time, reserved)} проверяет доступность клетки с учётом как статических препятствий, так и динамических резерваций (множество троек $(x, y, t)$, зарезервированных другими агентами).

\subsubsection{Класс Agent}

Класс \texttt{Agent} хранит текущую позицию агента $(x, y)$, целевую позицию $(goal\_x, goal\_y)$, построенный путь и индекс текущего шага. Метод \texttt{set\_path} назначает агенту новый путь, \texttt{next\_step} возвращает следующую позицию и продвигает индекс. Каждому агенту автоматически назначается уникальный цвет для визуализации.

\subsubsection{Класс Task}

Класс \texttt{Task} описывает задачу для MRTA-планировщика: уникальный идентификатор, целевая позиция $(x, y)$, приоритет, ссылка на назначенного агента и флаг завершения. Задачи создаются из целевых позиций сценария при выборе MRTA-режима.

\subsubsection{Класс SceneData}

Класс \texttt{SceneData} --- алгоритмо-независимое описание сценария: размеры сетки, позиции препятствий, стартовые позиции агентов и целевые позиции. Один и тот же сценарий может использоваться как для MAPF (цели --- конечные точки агентов), так и для MRTA (цели --- задачи для распределения). Сериализация выполняется в формат JSON, что обеспечивает переносимость и возможность повторного использования сценариев.

\subsection{Система регистрации алгоритмов}

Для обеспечения расширяемости системы реализован центральный реестр алгоритмов (модуль \texttt{registry}). Реестр представляет собой словарь, содержащий экземпляры всех доступных планировщиков, сгруппированные по типу (\texttt{"MAPF"} и \texttt{"MRTA"}).

Функция \texttt{get\_planner(algo\_type, display\_name)} выполняет поиск планировщика по типу и отображаемому имени. Такая архитектура позволяет добавлять новые алгоритмы без модификации существующего кода: достаточно реализовать подкласс \texttt{MAPFPlanner} или \texttt{MRTAPlanner}, импортировать его в \texttt{registry.py} и добавить экземпляр в соответствующий список.

Реестр используется как интерактивным интерфейсом (для формирования списка алгоритмов на экране выбора), так и подсистемой пакетного тестирования (для перебора комбинаций алгоритмов).

\subsection{Симулятор: планирование и выполнение}

Класс \texttt{Simulator} является центральным координирующим компонентом, объединяющим планирование и пошаговое выполнение для обоих режимов работы.

\subsubsection{Этап планирования}

Метод \texttt{plan()} выполняет полный цикл планирования с замером времени и памяти (\texttt{tracemalloc}):

\begin{enumerate}[arabic]
    \item \textbf{В режиме MAPF:} вызывается \texttt{planner.plan(world, agents)}, полученные пути назначаются агентам.
    \item \textbf{В режиме MRTA:} процесс двухэтапный:
    \begin{itemize}
        \item MRTA-планировщик распределяет задачи: \texttt{planner.plan(world, agents, tasks)};
        \item навигационный MAPF-планировщик строит маршруты к первым задачам: \texttt{\_nav\_all(agents)}.
    \end{itemize}
\end{enumerate}

Метод \texttt{\_nav\_all} является ключевым связующим элементом между MRTA и MAPF. Если задан навигационный планировщик (\texttt{nav\_planner}), он вызывается для совместного построения бесконфликтных путей всех агентов. В противном случае используется последовательный fallback в стиле CA*: агенты планируются один за другим, каждый --- с учётом путей предыдущих через общую таблицу резерваций.

\subsubsection{Этап выполнения}

Метод \texttt{step()} выполняет один временной шаг симуляции:

\begin{enumerate}[arabic]
    \item каждый агент, имеющий невыполненный путь, перемещается на следующую позицию;
    \item обновляется метрика SOC (каждое перемещение увеличивает счётчик на 1);
    \item в режиме MRTA проверяется завершение задач: если агент достиг позиции назначенной задачи, задача помечается как завершённая;
    \item при завершении задачи в queue-режиме из очереди извлекается следующая задача, и маршрут достраивается методом \texttt{\_path\_to\_single};
    \item при завершении задачи в online-режиме вызывается повторное планирование для освободившегося агента.
\end{enumerate}

\subsubsection{Сбор метрик}

Симулятор собирает шесть метрик, разделяя время планирования на составляющие:

\begin{table}
    \caption{Метрики, собираемые симулятором}
    \begin{tabularx}{\textwidth}{|l|X|}\hline
    Метрика              & Описание \\ \hline
    \texttt{makespan}    & Общее число шагов до завершения всех задач \\ \hline
    \texttt{soc}         & Sum of Costs --- суммарное число перемещений всех агентов \\ \hline
    \texttt{computation\_s} & Суммарное время планирования (секунды) \\ \hline
    \texttt{plan\_mrta\_s}  & Время, затраченное на MRTA-назначение \\ \hline
    \texttt{plan\_mapf\_s}  & Время, затраченное на навигационное планирование \\ \hline
    \texttt{memory\_peak\_mb} & Пиковое потребление памяти (МБ, через \texttt{tracemalloc}) \\ \hline
    \end{tabularx}
    \label{tab:metrics}
\end{table}

\subsection{Подсистема пакетного тестирования}

Подсистема пакетного тестирования обеспечивает автоматизированное проведение серий экспериментов с различными комбинациями алгоритмов и параметров среды.

\subsubsection{Конфигурация эксперимента}

Класс \texttt{BatchConfig} определяет все параметры пакетного эксперимента:

\begin{itemize}
    \item списки выбранных MAPF- и MRTA-алгоритмов;
    \item диапазон числа роботов: минимум, максимум, шаг;
    \item параметры карты: размер сетки, плотность препятствий, стратегия размещения;
    \item число сценариев на каждое значение числа роботов (для статистической достоверности);
    \item таймаут на один прогон;
    \item путь к директории вывода.
\end{itemize}

\subsubsection{Генерация сценариев}

Класс \texttt{SceneGenerator} создаёт сценарии для экспериментов. Поддерживаются четыре стратегии размещения агентов и задач:

\begin{itemize}
    \item \textbf{uniform} --- агенты и задачи размещаются на случайных свободных клетках равномерно;
    \item \textbf{clustered} --- агенты концентрируются в центре карты, задачи --- у периметра;
    \item \textbf{counter} --- агенты слева, задачи справа (встречные потоки);
    \item \textbf{min\_dist} --- гарантируется минимальное расстояние между стартовыми позициями.
\end{itemize}

Генератор выполняет проверку достижимости каждой пары старт--цель с помощью BFS и повторяет генерацию при обнаружении недостижимых целей (до 20 попыток). Также поддерживается импорт готовых JSON-карт.

\subsubsection{Запуск и вывод результатов}

Класс \texttt{BatchRunner} перебирает все комбинации (MRTA-алгоритм $\times$ MAPF-алгоритм $\times$ число роботов $\times$ seed сценария) и для каждой создаёт экземпляр симулятора. Таймаут реализован через \texttt{threading.Timer}: при превышении установленного времени прогон помечается как неуспешный.

Результаты агрегируются в \texttt{pandas.DataFrame} и экспортируются модулем \texttt{plotter}:

\begin{itemize}
    \item CSV-файл со всеми метриками каждого прогона;
    \item четыре графика (Matplotlib): Makespan, SOC, время планирования и Success Rate в зависимости от числа роботов.
\end{itemize}

Вывод сохраняется в папку с временной меткой вида \texttt{batch\_results/YYYYMMDD\_HHMMSS/}.
